\section{Experiments and Results}
\label{sec:results}

\begin{table}[t]
\centering
\footnotesize
\caption{The audit. Every row is the \emph{published} code, run with the authors' own script and configuration in an environment pinned to their own requirements; no method is re-implemented. \emph{Pub.\ I} and \emph{II} are the two splits the SpikingGCN paper itself reports, I being 20/500/1000 and II an 8:2 random split; where a paper reports one split, only that column is filled. \emph{Reproduced} is what the released code gives, averaged over $n$ seeds (multi-seed extensions of the single-run cells are reported in the text). The fourth row is SpikeGraphormer with one flag changed to the public split and everything else fixed.}\label{tab:audit}
\begin{tabular}{llrrrrr}
\toprule
method & dataset & pub.\ I & pub.\ II & reproduced & $\Delta$ & $n$ \\
\midrule
SpikingGCN & Cora & 80.7 & 88.7 & 50.0 & $-30.7$ & 5 \\
 & PubMed & 77.6 & 89.3 & 64.3 & $-13.3$ & 3 \\
 & CiteSeer & 72.5 & 77.6 & 70.6 & $-1.9$ & 1 \\
SpikeGraphormer & Cora & $84.8{\pm}0.7$ & --- & 84.7 & $-0.1$ & 3 \\
 & Cora, public split & --- & --- & 74.0 & $-10.8$ & 1 \\
SpikeGCL & Cora & $87.4{\pm}0.6$ & --- & 86.6 & $-0.8$ & 1 \\
 & CiteSeer & $77.6{\pm}0.6$ & --- & 77.3 & $-0.3$ & 1 \\
\bottomrule
\end{tabular}
\end{table}


\tabref{tab:audit} is the audit: three pipelines, seven rows, and three kinds of disagreement with the
literature. We take them in turn, because the remedy differs in each case.

\subsection{SpikingGCN: a dataset-dependent gap}
\label{sec:sgcn}

Run on the configurations its \texttt{models\_conf.json} ships, the released training script gives
$50.0$ on Cora over five seeds against a reported $80.7$, $64.3$ on PubMed over three seeds against
$77.6$, and $70.6$ on CiteSeer on one seed against $72.5$. The CiteSeer cell reproduces to within two
points; the Cora cell is thirty points short.

That one dataset reproduces and another does not is the most informative part of the result. It rules
out the mundane explanations --- that we could not execute the code, that the pinned environment is
wrong, or that the loader is misconfigured --- because the same script, in the same environment, with
the same loader reaches the published number on CiteSeer. The failure is specific to Cora and, to a
lesser degree, PubMed. A reader who took the Cora result alone could reasonably have concluded that
our reproduction was broken; with all three cells in \tabref{tab:audit}, that conclusion is not
available.

We therefore concentrated the diagnosis on Cora. The model is a two-hop polynomial propagation
followed by one spiking layer, and its input reaches the encoder through a Poisson sampler that
reads each entry as a firing probability. Several input conventions are plausible, and we
re-measured five of them through the audited pipeline, one seed each, with the transform inserted at
the loader's feature stage:

\begin{itemize}\itemsep2pt
\item \emph{Raw binary bag-of-words, self-loops} (faithful; what the library actually returns,
verified by dumping the array rather than by reading its source): $51.5$.
\item \emph{Raw binary bag-of-words, no self-loops}: $51.5$ --- the self-loop choice is not
decisive.
\item \emph{Row-normalised features} --- what a reader might reasonably assume: $13.0$, i.e.\ no
learning at all, because the entries are then of order $10^{-4}$ and the encoder never fires.
\item \emph{Row-normalised features, then encoder input rescaled to $[0,1]$ by its maximum} --- the
conventional fix for a Poisson encoder: also $13.0$. The rescaling restores the range but not the
per-node scale the sampler needs.
\item \emph{Row-normalised, no self-loops}: $13.0$ --- the cause is the normalisation, not the
self-loops.
\end{itemize}

Row-normalisation is the decisive cause; every variant that carries it collapses to $13.0$ and
every variant without it holds $51.5$. We list the failures because they are what a reader will
produce if they guess.

What we cannot exclude is that the published Cora number came from a state of the code, or a
per-dataset configuration, that was never released. The repository carries no commit tag
corresponding to the paper, no checkpoint and no container image, so the gap is unexplained rather
than attributed. The dataset dependence constrains any explanation: it has to account for why the
same script reproduces on CiteSeer while missing by thirty points on Cora. Our claim is the narrow
one --- the released code, in the pinned environment, does not reach the published Cora number under
any of the five measured input variants --- not that the authors made an error. A container image or
a commit hash would settle the question.

\subsection{SpikeGraphormer: a reproduction, and a split worth ten points}
\label{sec:split}

\begin{table}[t]
\centering
\footnotesize
\caption{SpikeGraphormer, Cora, everything fixed but the split. The only difference between the rows is one command-line flag; the code, the hyper-parameters, the seed and the number of epochs are identical. \emph{Final} is the accuracy at the epoch selected by validation, \emph{best} the best test accuracy over training. The published $84.8\pm0.7$ matches the first row.}\label{tab:split}
\begin{tabular}{lrr}
\toprule
split & final & best \\
\midrule
random, 20 per class (their \texttt{run.sh}) & 84.7 & 85.6 \\
public Planetoid (\texttt{--protocol semi}) & 74.0 & 75.9 \\
\bottomrule
\end{tabular}
\end{table}


The second pipeline behaves in the opposite way. Run with the arguments in its own \texttt{run.sh}, it
reaches $84.7$ over three runs against a published $84.8\pm0.7$: a difference of one tenth of a
point, well inside the reported standard deviation. If the audit stopped here it would report a
clean reproduction.

It does not stop here, because that configuration uses a random class-balanced split, and the paper's
footnote defers the split to the code. \tabref{tab:split} changes one flag --- from a random
class-balanced split to the fixed-split path --- and everything else stays identical: same code, same
hyper-parameters, same seed, same epochs. The result is $74.0$ instead of $84.7$.

Ten points is larger than the margin between most of the methods the field compares. A reader who
takes $84.8$ from the paper and sets it beside a public-split baseline is not comparing two methods;
they are comparing two tasks. Nothing in the paper says so, and nothing in the code says so either
--- the flag is buried in a shell script among many.

\subsection{SpikeGCL: two reproductions}
\label{sec:gcl}

The third pipeline reproduces on both datasets we ran, and that is what makes the audit
interpretable: the disagreements above are not our inability to run code.

On Cora, with the command in its README, the released code reaches $86.6$ against a published
$87.4\pm0.6$; on Citeseer it reaches $77.3$ against $77.6\pm0.6$. The Citeseer difference
($-0.3$ points) sits inside the published $\pm 0.6$ spread; the Cora difference ($-0.8$ points)
does not. Re-running the pipeline at four seeds gives a measured spread of $0.13$ points on Cora
and $1.24$ points on Citeseer, and a same-nominal-seed re-run of the Citeseer cell differed by
$1.0$ points from the originally audited run --- recorded as part of the sensitivity evidence.
The comparison against the published $\pm 0.6$ is the only spread
basis used here; a spread measured on SpikingGCN's Cora seeds ($0.48$--$0.53$) belongs to that
pipeline and not to SpikeGCL, and no PubMed range exists for SpikeGCL because its paper reports no
PubMed cell.

\subsection{One comparison source, two splits, no label}
\label{sec:mixed}

The last finding is not about a pipeline we ran but about the tables we read while choosing what to
run, and it explains the discrepancy that motivated the audit.

SpikingGCN's paper prints \emph{both} of its splits for every method it compares, and defines them
in its own words: ``In this split method (i.e., Split I), 20 instances from each class are sampled
as the training datasets'' and ``For the second data split (i.e., Split II), the ratio of training
to testing is 8:2, and 20\% of training samples is further used for validation''
\citep{zhu2022spiking}. The 500 validation and 1000 test nodes of Split~I are the standard
Planetoid sizes their loader leaves untouched, and the sampling is visible in their own
\texttt{sample\_per\_class(20, $\cdot$)} call. The two splits are not close. For SpikingGCN itself
the difference is $8.0$ points on Cora, $5.1$ on Citeseer and $11.7$ on PubMed --- larger than the
spread between the top methods in any of the field's comparison tables.

Against that, the two later papers we audited draw baseline rows from the two different splits of
the same comparison source, and the
inference is anchored on the quoted definitions rather than on paraphrase. SpikeGraphormer's split
is defined in code (their \texttt{run.sh} sets \texttt{-{-}rand\_split\_class -{-}valid\_num 500
-{-}test\_num 1000} for the published setting and \texttt{-{-}protocol semi} for the public split);
its table lists GCN at $81.5$ on Cora, which is consistent with the Split~I family.
SpikeGCL's table lists SpikingGCN at $87.7$ on Cora, $77.7$ on Citeseer and $87.5$ on PubMed ---
consistent, to within two points, with SpikingGCN's \emph{Split~II} numbers ($88.7$, $77.6$, $89.3$),
not its Split~I numbers ($80.7$, $72.5$, $77.6$). We phrase both readings as ``consistent with'':
the tables do not label the split, and the point is the missing label.

Neither table says which split it is quoting. A reader assembling a baseline table from these two
papers therefore mixes the splits without being told, and the resulting table carries a systematic
error of five to twelve points on some rows and not others. This is not a citation mistake in the
ordinary sense --- both numbers are correct and both papers report what they measured --- and that is
precisely why it has survived: there is no error to find, only an unlabelled choice.


\subsection{Split gaps against measured seed noise}
\label{sec:transport}

The split gaps are large in points, but a gap matters only against the noise a single pipeline
actually exhibits. We therefore measured seed spreads for the pipelines themselves: four seeds of
SpikeGCL give $0.13$ points on Cora and $1.24$ on CiteSeer; SpikingGCN's five Cora seeds give a
s.d.\ of $1.97$ points (per-seed $48.1$--$52.6$), and its three PubMed runs a s.d.\ of $5.37$
points; one same-nominal-seed CiteSeer re-run of SpikeGCL differed by $1.0$ points from the
audited cell, which the sensitivity analysis retains.

Matching each split gap to the noise of the pipeline it belongs to, the picture is consistent but
modest. Cora's $8.0$-point gap is $4.06$ times SpikingGCN's own seed s.d.; CiteSeer's $5.1$ is
$4.11$ times a proxied noise (the matching SpikingGCN cell is single-seed, so the SpikeGCL spread
is used and labelled as such); PubMed's $11.7$ is $2.18$ times its s.d. Against our three-sigma
bar --- fixed in the audit protocol before the multi-seed runs, as the pre-registration rule for
this comparison --- the verdicts are marginal ($1.35\times$, $1.37\times$) on Cora and CiteSeer and
fail on PubMed ($0.73\times$), where split choice and seed noise are the same order and cannot be
separated at the measured seed count. There is no dataset where the split gap dwarfs noise by an
order of magnitude; there is also none where the gap is inside noise except PubMed's confounded
cell. The transportability claim is therefore made per dataset at these magnitudes, not as a
blanket effect.

\subsection{The population, not a sample}
\label{sec:census}

\begin{table}[t]
\centering
\footnotesize
\caption{The population, not a sample. \emph{Declares} means the arXiv record or the paper itself gives a repository URL; the four that do not cannot be audited by anyone, and we checked each record before writing this column. The audit is therefore complete over the auditable population as of the date of writing.}\label{tab:census}
\begin{tabular}{llll}
\toprule
method & year & declares code & status \\
\midrule
SpikingGCN \citep{zhu2022spiking} & 2022 & yes & audited \\
SpikeGCL \citep{li2023bit1} & 2024 & yes & audited \\
SpikeGraphormer \citep{sun2024spikegraphormer} & 2024 & yes & audited \\
SiGNN \citep{chen2024signn} & 2024 & no & not auditable \\
DGSNN \citep{yin2024dgsnn} & 2024 & no & not auditable \\
MSG \citep{sun2024riemannian} & 2024 & no & not auditable \\
GSG \citep{zhang2025geometry} & 2025 & no & not auditable \\
\bottomrule
\end{tabular}
\end{table}


A three-method audit would be a sample if more methods were available. They are not.
\tabref{tab:census} lists the seven spiking-GNN methods in our related work and whether each declares
a code repository. Three do; we audited all three. The remaining four publish no code in their arXiv
record or their paper, so they cannot be audited by us or by anyone else, and any number attributed
to them in a comparison table is unverifiable in principle.

This changes what the audit is. It is not "three methods out of many"; within the seven-method
survey that frames this paper --- the spiking-GNN methods collected in our related work from the
2021--2026 literature, with repository status checked against each arXiv record and paper in
November 2026 --- the three auditable methods are the complete auditable set of that survey at the
time of writing. The claim is deliberately bounded: it does not assert a field-wide census, and the
reason the auditable set is small is itself part of the finding.

\subsection{What can now be placed side by side}

With three pipelines measured on the public split, four numbers become commensurable on Cora:
$50.0$ (SpikingGCN, its own code), $74.0$ (SpikeGraphormer, its own code with the split changed),
$86.6$ (SpikeGCL, self-supervised, its own code) --- and, shown for context and not counted as an
audited row, $80.5$ from the conversion study of our companion manuscript \citep{companion}, whose
own registered result files carry the value and whose authors' scripts are not part of this audit.
The published ordering --- SpikingGCN $80.7$, then SpikeGraphormer
$84.8$, then SpikeGCL $87.4$ --- is not the measured ordering, and the largest single move belongs to
the method that reproduces most exactly.

The four numbers measure three training paradigms on one dataset, so a ranking across paradigms
would be a category error. What they establish is narrower: a baseline table assembled from the
literature would have been wrong about all four.
